\section{Introduction}
\label{sec:introduction}

The equator is the least-volume closed minimal hypersurface of a unit
round sphere. The next volume level is a subtler question: minimality
alone gives no immediate control of topology, curvature, or singular
limits of a sequence of hypersurfaces. The natural comparison examples
are the minimal Clifford products
\[
\mathcal C_{k,m-k}
=S^k\!\left(\sqrt{\frac{k}{m}}\right)
\times S^{m-k}\!\left(\sqrt{\frac{m-k}{m}}\right)
\subset S^{m+1},\qquad 1\le k<m.
\]
The two factors have opposite principal curvatures, whose weighted sum
vanishes; these products are therefore minimal. For each integer
\(m\ge2\), set
\[
a_m=\min_{1\le k<m}\Vol_m(\mathcal C_{k,m-k}).
\]
The assertion that \(a_m\) is the first volume level above the equator
is the least-volume question in Yau's collection
\cite[Problem~31]{Yau1994}. The modern formulation in
\cite[Conjecture~1.12(ii)]{GeLi2021} calls it the Solomon--Yau conjecture.
We prove this formulation with its full immersion scope.

\begin{theorem}\label{thm:volume}
Let \(m\ge2\), let \(M^m\) be a closed connected smooth manifold, and let
\(F:M\to S^{m+1}\) be a smooth minimal immersion into the unit round sphere.
If \(F(M)\) is not totally geodesic, then
\[
\Vol\bigl(M,F^*g_{\mathrm{round}}\bigr)\ge a_m.
\]
\end{theorem}

Volume is measured on the domain, including covering multiplicities.
The theorem resolves the stated least-volume conjecture positively.
The threshold is attained by a Clifford product minimizing \(a_m\).

\subsection*{Earlier results and the geometric obstacle}

Simons' curvature identities underpin rigidity for minimal hypersurfaces
with parallel second fundamental form \cite{Simons1968}, explaining the
special role of the Clifford products. Cheng, Li, and Yau obtained a
dimension-dependent volume gap above the equator by heat-kernel
comparison \cite[Corollary~5]{ChengLiYau1984}; identifying the sharp
next level is a separate question. In dimension two, the sharp
area bound \(a_2=2\pi^2\), together with Clifford-torus rigidity, follows
from the Willmore theorem of Marques and Neves
\cite[Theorem~B and Remark~1.3]{MarquesNeves2014}.
Their canonical-family and min--max arguments connect an area comparison
to the topology of a family of surfaces. The present proof follows this
variational strategy, but the family and its topological test must work
in arbitrary hypersurface dimension.

Higher-dimensional progress includes the rotational results of Cheng,
Wei, and Zeng \cite[Theorem~1.1]{ChengWeiZeng2019}, and Viana's sharp
Clifford bound when both connected complementary regions are
antipodally invariant \cite[Theorem~1.4]{Viana2023}.
Related spectral progress includes Tang and Yan's proof that the first
positive Laplacian eigenvalue of a closed minimal isoparametric
hypersurface in the unit sphere equals its dimension
\cite[Theorem~1.1]{TangYan2013}.
Bernstein and Wang obtain near-sharp density bounds for regular minimal
cones of dimension at most six, under topological conditions on a
complementary component \cite[Theorems~1.1--1.2]{BernsteinWang2025}.
Here cone dimension is one greater than hypersurface dimension: the
density of the cone over a minimal hypersurface \(\Sigma^m\) is
\(|\Sigma|/|S^m|\). These density estimates indicate why the geometry of
cones is relevant to the volume question.

The nonembedded case has a useful independent reduction. A point with
\(q\) preimages forces domain volume at least \(q|S^m|\), as proved by
Nguyen \cite[Corollary~4.8 and Remark~4.9]{Nguyen2023} and by
Ge and Li \cite[Theorem~1.2 and Corollary~1.3]{GeLi2022}.
Since \(a_m<2|S^m|\), any counterexample to Theorem~\ref{thm:volume}
would be embedded. We give the short cone-monotonicity proof in
Proposition~\ref{prop:immersion-reduction}, retaining tangent and
coincident sheets. The more recent general embedded-hypersurface bounds of Ge and Li
\cite[Theorems~1.4 and 1.6]{GeLi2026II} remain below the sharp Clifford
threshold, which their Conjecture~1.3 still poses as the target.

After this reduction, the central difficulty is to minimize area among
non-equatorial counterexamples and then decrease that minimum.
Ordinary varifold compactness may produce singular limits, while a
variational construction must retain enough topology to avoid collapsing
to an equator. We address these two requirements together: a detecting
family gives a uniform gap above the equator, and a dimension induction
rules out singularities below the Clifford threshold.

\subsection*{Curvature, index, and the variational family}

The first construction controls both signs of the principal curvatures.
For a closed embedded non-equatorial minimal hypersurface \(\Sigma^m\),
with second fundamental form \(A\), a two-point maximum compares the
normal at one point with every other point of \(\Sigma\).
The reflection calculation is related to the two-point methods in
\cite[Section~2]{Brendle2013} and \cite[Section~2]{AHLW2015}.
We prove the required static differential
inequality directly, including diagonal contacts and repeated principal
curvatures. A smooth upper-test comparison then gives
\[
\int_\Sigma |A|^2\ge m|\Sigma|.
\]
This proves the lower-bound assertion of Perdomo's average-curvature
conjecture \cite[Conjecture~1.5]{GeLi2021}.
Combining it with the conformal and spectral mechanism of
\cite[Theorem~3.1]{Perdomo2004}, and keeping track of equality in each
spectral inequality, shows that index \(m+3\) forces a Clifford product
(Proposition~\ref{prop:index-rigidity}). A counterexample must
therefore have a further negative direction for area.

The average-curvature inequality and this embedded index classification
also appear in Nifa's recent preprint
\cite[Theorem~1.1 and Corollary~7.3]{Nifa2026}, which proves a stronger
Schr\"odinger quadratic-form inequality. His construction uses the sum
of two tangent-ball curvatures and a weak supersolution. Here a two-sided
maximum and a smooth Poisson comparison suffice for the integrated
bound. We give the curvature and index arguments in full, independently
of the volume comparison that follows.

The second construction turns that direction into a strict area
improvement. Form the Euclidean cone over \(\Sigma\), take its signed
distance function, and regard the graph as a hypersurface in a space
with one time coordinate. Lorentz boosts preserve the difference
between squared spatial length and squared time. Intersecting the
transformed graph with the unit spatial sphere produces level sets
whose areas are controlled on the fixed source \(\Sigma\).
At each closest-point contact, minimality makes a nonnegative Jacobian
matrix have fixed trace; its determinant bounds the area of the level.
This estimate works even when closest points are not unique.

The boost parameters and the level parameter give a family on
\(\mathbb{RP}^{m+3}\). Its boundary cycles have area at most
\(|\Sigma|\), and its fillings change to their complements around the
projective generator. The associated degree-one cohomology class has
nonzero \((m+3)\)-rd power. This is the topological detection used in
Section~\ref{sec:minmax}; it replaces a topology assumption on
\(\Sigma\) itself. The original hypersurface is the unique area
maximum, with a uniform gap away from its central parameter.
Its extra index direction removes that maximum while preserving the
detection and the absence of mass concentration
(Proposition~\ref{prop:lorentz-family}).

\subsection*{Compactness and the final induction}

Detection forces every such family to contain a boundary of a
half-volume region with zero centroid. Spherical isoperimetry and its
hemisphere equality case give a uniform area threshold
\(\gamma_m>|S^m|\) for these boundaries
(Lemma~\ref{lem:detection-gap}). The min--max theorem therefore produces
stationary integral varifolds above the equator, even after the
strict deformation.

The volume proof proceeds by induction on \(m\). Assuming only the
lower-dimensional volume statements, Section~\ref{sec:cones} proves
regularity for the subthreshold limits arising both from minimizing
sequences and from min--max. The mechanism is radial splitting:
tangents at nonzero points of a Euclidean cone split off a line,
reducing their classification to a smaller cone dimension. This gives
smooth connected links of multiplicity one without applying the
unknown dimension-\(m\) volume theorem to those links.

Compactness therefore attains a least counterexample of area \(H\),
with \(\gamma_m\le H<a_m\). Its index exceeds \(m+3\), so the deformed
family has maximum area strictly below \(H\). Min--max produces a
varifold of mass in \([\gamma_m,H)\), and the same regularity argument
makes it a smooth connected multiplicity-one hypersurface. Its mass
is above the equator and below \(H\): it is a smaller counterexample,
the desired contradiction. The induction starts at \(m=2\), with
no lower-dimensional volume hypothesis.

Section~\ref{sec:preliminaries} establishes the volume normalization and
immersion reduction. Section~\ref{sec:index} proves curvature and index
rigidity. Section~\ref{sec:minmax} establishes detection and the precise
min--max interface before Section~\ref{sec:family} constructs and deforms
the Lorentz-distance family. Section~\ref{sec:cones} proves the cone
induction and exact attainment, and Section~\ref{sec:conclusion}
completes the proof.
